\begin{abstract}
Unsupervised domain adaptation for panoramic semantic segmentation requires transferring from labeled perspective source domains to unlabeled panoramic target domains, facing challenges including projection distortion, cross-field-of-view dependencies, and pseudo-label noise. OmniSAM employs a SAM2-style memory mechanism to model panoramic view sequences, but interactions with numerous historical memory tokens significantly increase training costs. This paper proposes a memory-efficient training scheme consisting of Memory Token Compression (MTC) and Sparse Memory Retrieval (SMR): MTC aggregates spatially adjacent redundant memories, while SMR sparsely selects historical tokens by importance, thereby reducing the representation scale entering memory attention; pseudo-label quality optimization serves only as auxiliary exploration. In Cityscapes $\rightarrow$ DensePASS experiments with sam2\_s on 13 classes using a single seed, training time decreases by approximately 22.4\%, with mIoU changing from 60.26 to 60.16 and BFScore from 39.97 to 39.86. Ablation results further demonstrate that token budget reduction is directly correlated with attention overhead reduction, though stronger compression may affect boundaries and small objects. Qualitative comparisons show that both configurations maintain similar major scene structures, with local differences consistent with slight boundary metric degradation. These results support trading moderate accuracy for training efficiency, though cross-dataset and multi-seed stability remain to be verified.
\end{abstract}

\noindent\textbf{Keywords:} Panoramic semantic segmentation; Unsupervised domain adaptation; OmniSAM; Memory mechanism; Token compression; Efficient training
